% titel: Blatt 3, A9: Enthalpieänderungen von Ethanol
% skript: Mortimer Kap. 5 (Kalorimetrie, Verdampfung, Verbrennung, Bildung)

% Teil I: Werkzeuge, kein Lösungsweg.

Wird bei bekanntem $c_p$, Verdampfungs-, Verbrennungs- oder
Bildungsenthalpie eine Stoffmenge erwärmt, verdampft, verbrannt oder gebildet,
so berechnet man die jeweilige Enthalpieänderung.

\begin{werkzeugkasten}[Die vier Enthalpiebeiträge]
Für eine Stoffmenge $n = m/M$ (mit Masse $m$ und Molmasse $M$):
\begin{itemize}
  \item \textbf{Erwärmen} (Phasenzustand bleibt flüssig):
        $\Delta H = m\cdot c_p\cdot \Delta T$.
  \item \textbf{Verdampfen}: $\Delta H = n\cdot \Delta_{vap}H^\circ$.
  \item \textbf{Verbrennen}: $\Delta H = n\cdot \Delta_cH^\circ$.
  \item \textbf{Bilden aus den Elementen}: $\Delta H = n\cdot \Delta_fH^\circ$.
\end{itemize}
Vorzeichen zählen: Verbrennungs- und Bildungsenthalpie von Verbindungen sind
typischerweise negativ, die Verdampfungsenthalpie positiv.

\textit{Skript:} Mortimer Kap. 5 (Phasenumwandlung, Verbrennung, Bildung),
Übung 3/9.
\end{werkzeugkasten}

\begin{achtung}[Darauf achten]
\begin{itemize}
  \item Für Erwärmen die \emph{Masse} $m$ (in g) mit $c_p$ (in J\,g$^{-1}$K$^{-1}$)
        einsetzen; für Verdampfen/Verbrennen/Bilden die \emph{Stoffmenge}
        $n$ mit molaren Grössen (kJ/mol).
  \item Anzahl der Mol $n = 10\ \mathrm{g} / 46.07\
        \mathrm{g\,mol^{-1}}$ nicht vergessen.
  \item Vorzeichen: $\Delta_{vap}H > 0$ (endotherm), $\Delta_cH < 0$,
        $\Delta_fH(\text{flüssig}) < 0$.
  \item Unterschied konkret: $c_p$ bezieht sich auf Gramm, molare Enthalpien
        auf Mol --- nie beide gleichsetzen.
\end{itemize}
\end{achtung}
